\section{Introduction}\label{Introduction}

\subsection{Exchange rates}\label{Exchange rates}

Since the fall of the Bretton Woods system of fixed exchange rates in 1976 and the establishment of floating exchange rates with the Jamaica Accords, the need for accurate forecasting of exchange rates has become increasingly important. The primary reason for this is the need for effective risk management, especially for businesses, investors, currency traders, governments, and banks. Knowing whether the value of a currency will rise or fall ahead of time allows not only for the mitigation of adverse currency movements, but also for well-timed and potentially highly profitable investment decisions.

Unfortunately, exchange rates are notoriously difficult to predict as they exhibit a complex non-linear nature that is dependant on a wide range of variables including macroeconomic and microeconomic factors.

Given the challenges in forecasting exchange rates using traditional econometric models, researchers have increasingly turned to the application of artificial intelligence (AI) algorithms as a potential solution. AI-based approaches, such as neural networks, genetic algorithms, and machine learning techniques, have shown promise in their ability to capture the complex, nonlinear relationships inherent in exchange rate dynamics. These methods have the potential to outperform traditional forecasting models by identifying patterns and relationships that may be difficult for human analysts to discern.

\subsection{Literature review}\label{Literature review}


